\section{Opening the model}
\label{sec:method}
% claims: C02 C11 C12 C13 C14 C17 C19 C35 C38 C43
The interpretability method remains part of the work. Its purpose is to connect
an internal change to a measured output effect, while checking what the edit
actually changed. Its historical empirical claims require the repairs listed
later; the procedure itself does not depend on a favorable result.

\subsection{Adapt the model, then hold it fixed}
We adapt Qwen3-8B on non-insider user-days using QLoRA
\citep{qwen3,hu2022lora,dettmers2023qlora}. The base weights are frozen and
quantized; small low-rank matrices supply the learned updates. For a projection
matrix, the adapted weight is $W'=W+(\alpha/r)BA$. This makes it practical to
compare the same model with adapters enabled and disabled. The model is frozen
while we audit it. QLoRA supplies an accessible adaptation difference; it does
not establish that an LM is the best detector for these data.

\subsection{Isolate the adaptation difference}
At token position $t$ and layer $\ell$, compare the adapted and base states:
\begin{equation}
 \delta_t=h^{(\ell),\mathrm{ada}}_t-h^{(\ell),\mathrm{base}}_t.
\end{equation}
This is the representation difference induced by adaptation on that input.
It is not automatically a unique causal decomposition of the network.

\subsection{Find a sparse approximation}
Standardize each coordinate of $\delta_t$ and fit a sparse autoencoder (SAE):
\begin{equation}
 q_t=(\delta_t-\mu)\oslash\sigma,\qquad
 z_t=\operatorname{TopK}(\operatorname{ReLU}(Eq_t+b_e)),\qquad
 \widehat q_t=Dz_t+b_d.
\end{equation}
Only a few coefficients of $z_t$ are retained. The dictionary columns of $D$
are candidate directions of change, learned to reconstruct the observed deltas
\citep{cunningham2023sparse,gao2024scaling}. Sparse approximation is not a proof
that these directions are unique concepts. Any cross-dictionary comparison must
account for normalization: physical directions are $\sigma\odot D_{:,f}$.

\subsection{Select candidates, edit, and measure}
An activation difference between incident-window and reference examples can
rank candidate features. Selection is a hypothesis-generating step: labels,
reference-user eligibility, and the dictionary sample all affect the answer.
To test a selected set, edit its coefficients, continue the forward pass, and
measure the change in a declared output such as behavior-target NLL. Compare
against active features and random directions with the same edited support and
matched realized perturbation size. Keep a same-configuration unedited forward
pass as the reference. We measure both the requested edit and the change that
actually survives low-precision rounding.

A change in output establishes sensitivity to that specified edit. It does not
by itself prove that the features encode maliciousness, that the edit improves
anomaly ranking, or that the selected set is strictly necessary. Best-of-search
edits also differ from an intervention policy chosen before seeing its outcome.

\subsection{Locate activation without overinterpreting it}
For a fixed set of examples, activation mass at input class $c$ is
\begin{equation}
 m_f(c)=\frac{\sum_{x,t:\,\mathrm{class}(t)=c}z_{t,f}(x)}
                 {\sum_{x,t}z_{t,f}(x)}.
\end{equation}
This describes where a feature activates. A session-position hidden state can
use an earlier profile. High profile-position mass in selected features is
neither a complete account of the detector nor proof of a profile-only causal
mechanism. Corrected attribution must use the corrected selected features and
example population. The earlier profile-versus-behavior headline is pending
that recovery.
